\documentclass[11pt,a4paper]{article}
\usepackage[T1]{fontenc}
\usepackage{lmodern,amsmath,amssymb,booktabs,microtype}
\usepackage[margin=24mm]{geometry}
\usepackage[colorlinks=true,urlcolor=blue,linkcolor=blue]{hyperref}
\newcommand{\EE}{\mathbb E}
\newcommand{\tr}{\operatorname{tr}}
\newcommand{\NN}{\mathcal N}
\newcommand{\XX}{\mathscr X}
\newcommand{\DD}{\mathfrak D}
\title{Adjacent-charge propagation as Gaussian smoothing\\A finite-time transfer estimate with all contractions retained}
\author{Alan Day \and Chengheng\thanks{AI research assistant (OpenAI).}}
\date{Local research note --- 7 November 2026}
\begin{document}\maketitle
\begin{abstract}
The ordinary charge-removal propagator admits an approximation by
Gaussian Ornstein--Uhlenbeck smoothing of the propagator in the original
charge sector. Its annealed normalized Hilbert--Schmidt error is
$O(N^{-1})$ on compact real-time intervals, after the full infinite-temperature
charge average. We first compute both directional transfer defects,
including the nonzero reverse contraction source. Polar normalization
then gives a small creation defect; adjoints and Gaussian-chaos tails
extend it to the complete ordinary evolution. The preceding rank--chaos
bound converts the resulting rank filter to a disorder-space smoothing.
This is a propagation result; force insertions and their energy-regressed
fluctuations remain the next calculation.
\end{abstract}

\section{Model and principal statement}
Use standard CAR, $\hbar=1$, string time, and
$H=\sum_{I,J}J_{IJ}c_I^\dagger c_J$, with overlapping $p$-body index sets,
$J=J^\dagger$, and
\begin{equation}
 \EE J_{IJ}J_{KL}=w_N\delta_{IL}\delta_{JK},\quad
 w_N=4^p\frac{J_s^2}{\lambda_N}\left(\frac{p!}{N^p}\right)^2,
 \qquad \lambda_N=p^2/N.
\end{equation}
All charge sectors share the same $J$. Write $h=H_m$, $k=H_{m-1}$,
$\|X\|_{2,m}^2=\EE\tr_mX^\dagger X$, $\tr_m=\operatorname{Tr}/\binom Nm$,
and $b_m=2^{-N}\binom Nm$. Define
\begin{equation}
 T_m(X)=\frac1m\sum_i c_i^\dagger Xc_i,\qquad
 F_m(s)=T_m(e^{isk}),\qquad \alpha_m=(1-p/m)_+.
\end{equation}
For $0\le a\le1$, the disorder Ornstein--Uhlenbeck map is
\begin{equation}
 (\mathcal O_a f)(J)=\EE_{J'}f(aJ+\sqrt{1-a^2}J'),
 \qquad J'\text{ an independent copy}.
 \label{eq:ou}
\end{equation}
As $N\to\infty$, $\lambda_N$ stays in a compact subset of $(0,\infty)$,
and $J_s$ and $S$ are fixed. We prove
\begin{equation}
 \boxed{\quad
 \sup_{|s|\le S}\sum_{m=1}^N b_m
 \|F_m(s)-\mathcal O_{\alpha_m}(e^{ish})\|_{2,m}^2
 =O_S(N^{-2}).\quad}
 \label{eq:main}
\end{equation}
The approximation preserves the original sample $J$; only the
conditional average in \eqref{eq:ou} uses a new copy.

\newpage
\section{The channel spectrum and a forward defect}
For the finite identities assume $m>p\ge1$. Set
\begin{equation}
 n=N-m+1,\quad \alpha=(m-p)/m,\quad
 C=w_N\binom mp\binom{N-m+p}p,\quad
 C_-=C\alpha\frac{n+p}{n}.
\end{equation}
Subscripts $+$ and $-$ denote the upper and lower charge operator
spaces with their normalized Hilbert--Schmidt inner products. The
adjoint of $T=T_m$ is
\begin{equation}
 A=T^\dagger,\qquad A(Y)=\frac1n\sum_i c_iYc_i^\dagger.
\end{equation}
Let $P_r^\pm$ be the multiplicity-free neutral $U(N)$ rank projectors,
and $\XX_\pm=\sum_r x_rP_r^\pm$, $x_r=r(N+1-r)$. Then
\begin{equation}
 AT=1-\frac{\XX_-}{mn},\qquad
 TA=1-\frac{\XX_+}{mn},\qquad T\XX_-=\XX_+T.
 \label{eq:spectrum}
\end{equation}
These spectra follow from the known embedded-GUE covariance spectrum
\cite{Kota}: the one-body case gives $TA$ directly. CAR, or the
primitive-rank norm ratio, gives $AT$. A rank present on only one
side has $x_r=mn$ and is killed by the corresponding map. Both
$T$ and $A$ are unital, trace preserving in these normalized traces,
and contractions. The spectrum itself is an existing result's application.

Use an independent pair $(u,v)$ sharing one new Gaussian coupling
in the two charges. The earlier charged-covariance interface gives
\begin{equation}
 Tv=\alpha u,\qquad \EE(uc_i^\dagger v)=\alpha Cc_i^\dagger,
 \qquad \EE\,uT(vX)=\alpha C TX.
 \label{eq:cross}
\end{equation}
For deterministic lower-rank input $X_r$, the one-step Casimir mean
and \eqref{eq:spectrum} imply
\begin{align}
 \EE\|T(vX_r)\|_+^2
 &=\left(\alpha^2C-\frac{\kappa_1^-x_r}{mn}\right)\|X_r\|_-^2,
 &\kappa_1^-&=\frac{C(m-p)(m-p-1)}{m(m-1)}.
\end{align}
Subtracting the two cross terms using \eqref{eq:cross} proves
\begin{equation}
 \boxed{\ \EE\|T(vX_r)-\alpha uTX_r\|_+^2
 =\gamma x_r\|X_r\|_-^2,\qquad
 \gamma=\frac{Cp(m-p)}{m^3(m-1)n}.\ }
 \label{eq:forward}
\end{equation}
In real Gaussian coordinates $u=\sum_\beta g_\beta B_\beta^+$,
$v=\sum_\beta g_\beta B_\beta^-$, this is an exact coefficient
quadratic form. Its defining positive operator is equivariant, so
different input ranks have zero cross terms. If $X$ depends on the
original $J$, the new copy is averaged with $X$ held fixed.

\newpage
\section{The reverse source is explicit and orthogonal}
The reverse-adjoint direction contains a source even on the identity.
Put
\begin{equation}
 V=v-\alpha A(u)=\frac{\XX_-v}{mn},\qquad
 D_u(Y)=A(uY)-A(u)A(Y).
\end{equation}
There is an exact decomposition
\begin{equation}
 vA(Y)-\alpha A(uY)=V A(Y)-\alpha D_u(Y).
 \label{eq:reverse-split}
\end{equation}
The two terms are orthogonal in mean:
\begin{equation}
 \EE\langle V A(Y),D_u(Y)\rangle_-=0.
 \label{eq:orth}
\end{equation}
Indeed the equivariant map $Z\mapsto\EE[VZu]$ acts by one scalar
$j_V$ on the unique fundamental charged type $c_i$. Thus
$\EE[V A(uY)]=j_VA(Y)$ and $\EE[V A(u)]=j_VI$, proving
\eqref{eq:orth}. No independence of the two terms is used.

For clarity the remaining coefficient can also be computed without
new recoupling data. Let $c_+=e_+(1-\mu_1^+)$ be the upper
one-step Casimir mean from the identity, and put
\begin{equation}
 C_A=\EE\tr_-[A(u)]^2=C\left(1-\frac{c_+}{mn}\right).
\end{equation}
The fundamental charged covariance gives
$\EE[u c_i^\dagger A(u)]=C_Ac_i^\dagger$. Consequently
\begin{align}
 \EE\|D_u(Y_r)\|_-^2
 &=\left\{C\left[1-\frac{\mu_1^+x_r+c_+}{mn}\right]
      -C_A\left[1-\frac{x_r}{mn}\right]\right\}\|Y_r\|_+^2\notag\\
 &=\frac{Cp x_r}{m n^2(n+p-1)}\|Y_r\|_+^2.
\end{align}
The known second Casimir moment for $v$ gives
\begin{equation}
 \nu_0=\EE\tr_-V^2
 =\frac{Cp^2(m-p)N(N+1)}{m^3n(n+p-1)}.
\end{equation}
Equations \eqref{eq:reverse-split}--\eqref{eq:orth} therefore determine
the entire reverse quadratic form:
\begin{equation}
 \boxed{\quad
 \EE\|vA(Y_r)-\alpha A(uY_r)\|_-^2
 =\left[\nu_0\left(1-\frac{x_r}{mn}\right)
  +\frac{\alpha^2Cp x_r}{mn^2(n+p-1)}\right]\|Y_r\|_+^2.
 \quad}
\end{equation}
At central filling the source has norm of order $p/N$, whereas
the centered defect on ranks $r\le kp$ has squared norm
$O(C\lambda_N k/N^2)$. Simply copying the forward estimate to
annihilation would lose the source. Polar normalization handles
both directions while retaining its effect in the channel amplitude.

\newpage
\section{Polar normalization makes the dynamical defect small}
Take the polar decomposition on normalized operator Hilbert spaces,
\begin{equation}
 T=\mathcal U S_-=S_+\mathcal U,\qquad
 S_\pm=\sqrt{1-\XX_\pm/(mn)},\qquad
 s_r=\sqrt{\frac{(m-r)(n-r)}{mn}}.
\end{equation}
$\mathcal U$ is a deterministic, rank-preserving partial isometry
and $\mathcal U I=I$. We never invert $s_r$ at a zero mode.
In this and the next section take $N\ge16$, $p\le N/16$,
$N/4\le m\le3N/4$, and
\begin{equation}
 B_m=\lambda_N+\lambda_N^2+
       \lambda_N\frac{(m-n)^2}{N}.
\end{equation}
For $r\le kp$ and $(k+1)p\le N/8$, all ranks reached by one
multiplication are common to both charges, and $s_t\ge1/2$.
Writing $E_\beta=TB_\beta^- -\alpha B_\beta^+T$ as left
multiplication maps, the rank-$r$ to rank-$t$ identity is
\begin{equation}
 P_t^+(\mathcal U B_\beta^- -B_\beta^+\mathcal U)X_r
 =\frac{P_t^+E_\beta X_r
       +(\alpha s_r-s_t)P_t^+(B_\beta^+\mathcal U X_r)}{s_t}.
 \label{eq:polarstep}
\end{equation}
The product triangle gives $\delta=p+r-t\in[0,2p]$. The preceding
rank-chain result gives, for the upper multiplication law,
\begin{equation}
 \EE_r\delta^2\le10^6[\lambda_N(k+1)+\lambda_N^2(k+1)^2].
\end{equation}
On this interval $|s_t-s_{r+p}|\le32\delta/N$. Also, exactly,
\begin{equation}
 \alpha^2s_r^2-s_{r+p}^2
 =\frac{p}{mn}\left[\alpha(m-n)+\frac{2r(n-r)}m
                                      -\frac{p x_r}{m^2}\right],
\end{equation}
so that
\begin{equation}
 |\alpha s_r-s_{r+p}|
 \le\frac{32p|m-n|}{N^2}+\frac{640kp^2}{N^2}.
\end{equation}
Combine these estimates, \eqref{eq:forward}, and
$\gamma x_r\le4096C\lambda_N k/N^2$ in \eqref{eq:polarstep}.
Different input-rank cross terms vanish by equivariance. For Gaussian
creation $\mathcal C_\pm=\sum a_\beta^\dagger L_{B_\beta^\pm}$,
symmetrization of a degree-$k$ coefficient tensor costs only $k+1$:
\begin{equation}
 \|(\mathcal U\mathcal C_- -\mathcal C_+\mathcal U)\psi_k\|^2
 \le\frac{K_0 C B_m}{N^2}(k+1)^3\|\psi_k\|^2,
 \qquad K_0=10^{12}.
 \label{eq:creation}
\end{equation}
This constant is deliberately loose. The filling drift in $B_m$
must be retained; a uniform $O(N^{-1})$ norm claim throughout the
entire broad central interval would be incorrect.

\newpage
\section{All chaoses and the full real-time propagator}
Work in the Gaussian $L^2$ space of equivariant matrix functions;
$\NN$ is the chaos number. Degree $k$ has matrix rank at most $kp$.
Put $K=\lfloor N/(8p)\rfloor-1$. On degrees at most $K+1$, the
partial isometry is unitary between the complete equivariant
subspaces. If $E=\mathcal U\mathcal C_- -\mathcal C_+\mathcal U$,
the annihilation defect on these blocks is
\begin{equation}
 \mathcal U\mathcal C_-^\dagger-\mathcal C_+^\dagger\mathcal U
 =\mathcal U E^\dagger\mathcal U.
\end{equation}
Thus \eqref{eq:creation} controls both directions. Gaussian
creation and annihilation conventions are standard \cite{ChenLiu}.
Let $\mathcal H_\pm=L_{H_\pm}=\mathcal C_\pm+\mathcal C_\pm^\dagger$
and $G=\mathcal U\mathcal H_- -\mathcal H_+\mathcal U$.
With $M=\max_j C_j$,
\begin{align}
 \|G\psi_{\le K}\|^2
 &\le\frac{4K_0 M B_m}{N^2}\|(\NN+1)^{3/2}\psi_{\le K}\|^2,\\
 \|G\psi_{>K}\|^2
 &\le16M\|(\NN+1)^{1/2}\psi_{>K}\|^2.
\end{align}
The second line uses only $\|\mathcal U\|\le1$ and
$\|\mathcal H_\pm\psi\|^2\le4C_\pm\|(\NN+1)^{1/2}\psi\|^2$.

For $\psi(t)=e^{itk}$, the preceding ordinary-propagator proof gives
\begin{align}
 \sum_j\zeta^j\|\psi_j(t)\|^2&\le e^{C_-t^2(\zeta-1)},\quad\zeta\ge1,\\
 \|(\NN+1)^{3/2}\psi(t)\|^2&\le P_3(C_-t^2),
 &P_3(z)&=1+7z+6z^2+z^3.
\end{align}
Its factorial moments also give
$\|(\NN+1)^{1/2}\psi_{>K}\|^2
\le2^{-(K+1)}(1+2z)e^z$, $z=C_-t^2$.
For $|s|\le S$, define $z_S=MS^2$ and
\begin{equation}
 R_m(S)=S^2\left[
 \frac{8K_0 M B_m}{N^2}P_3(z_S)
 +32M2^{-(K+1)}(1+2z_S)e^{z_S}\right].
\end{equation}
The exact Hilbert-space Duhamel formula, using $\mathcal U I=I$, is
\begin{equation}
 \mathcal Ue^{isk}-e^{ish}
 =i\int_0^s e^{i(s-t)h}G(e^{itk})\,dt.
\end{equation}
Finite-$N$ Gaussian Sobolev domains justify the formula; the left
factor is unitary for real time. Hence
\begin{equation}
 \|\mathcal Ue^{isk}-e^{ish}\|_+^2\le R_m(S),\qquad
 \|F_m(s)-S_+e^{ish}\|_+^2\le R_m(S).
 \label{eq:rankproxy}
\end{equation}
For fixed $\lambda_N>0$ the chaos-tail term is exponentially small
in $\sqrt N$. No truncated time series is used.

\newpage
\section{From the physical-rank filter to disorder OU smoothing}
This step uses the full ordinary rank--chaos theorem from the
preceding stage. Define $\mathcal R=\sum_r rP_r^+$ and
$\DD=p\NN-\mathcal R\ge0$ on equivariant functions.
For every allowed upper rank, $u=x_r/(mn)\in[0,1]$, and
\begin{equation}
 0\le1-u/2-\sqrt{1-u}\le u^2/2.
\end{equation}
Using $m,n\ge N/4$, $x_r=r(m+n-r)$, and $N+1\le2N$, one obtains
\begin{equation}
 |s_r-(1-r/m)|\le
 \frac{8r|m-n|}{N^2}+\frac{520r^2}{N^2}.
\end{equation}
For integer $k\ge0$, Taylor's theorem on $0\le p/m\le1$ gives
$|\alpha^k-1+kp/m|\le k(k-1)p^2/(2m^2)$ (zero for $k=0,1$).
Consequently, for all $r\le kp$, without a low-rank restriction,
\begin{equation}
 \boxed{\quad
 |s_r-\alpha^k|\le\frac{4(kp-r)}N
 +\frac{8kp|m-n|}{N^2}+\frac{528k^2p^2}{N^2}.\quad}
 \label{eq:filterbound}
\end{equation}
Orthogonal rank and chaos decompositions then give
\begin{align}
 \|(S_+-\alpha^{\NN})e^{ish}\|_+
 &\le\frac4N\|\DD e^{ish}\|_+
 +\frac{8p|m-n|}{N^2}\|\NN e^{ish}\|_+\notag\\
 &\hspace{12mm}+\frac{528p^2}{N^2}\|\NN^2e^{ish}\|_+.
\end{align}
For $z=Cs^2$, the established bounds are
\begin{align}
 \|\DD e^{ish}\|_+^2&\le4\cdot10^6 C(\lambda_N+\lambda_N^2)s^2P_3(z),\\
 \|\NN e^{ish}\|_+^2&\le z+z^2,\\
 \|\NN^2e^{ish}\|_+^2&\le P_4(z),&
 P_4(z)&=z+7z^2+6z^3+z^4.
\end{align}
Thus a finite bound for the squared filter error is
\begin{align}
 Q_m(S)={}&\frac{192\cdot10^6 M(\lambda_N+\lambda_N^2)S^2}{N^2}P_3(z_S)\notag\\
 &+\frac{192\lambda_N(m-n)^2}{N^3}(z_S+z_S^2)
 +\frac{836352\lambda_N^2}{N^2}P_4(z_S).
\end{align}
Since $\alpha^{\NN}=\mathcal O_\alpha$, combining with
\eqref{eq:rankproxy} yields the central-sector bound
\begin{equation}
 \|F_m(s)-\mathcal O_\alpha(e^{ish})\|_+^2
 \le2R_m(S)+2Q_m(S).
 \label{eq:finite}
\end{equation}
This smoothing acts on the Gaussian coupling, not on the spectral
energy variable. Its parameter $\alpha$ is distinct from the
previous energy-Mehler parameter $q=e^{-4\lambda}$.

\newpage
\section{Full charge average and the next memory calculation}
The constants remain uniform in the full-charge ensemble. Indeed
\begin{equation}
 C_m\le\frac{J_s^2}{\lambda_N}
 \prod_{j=0}^{p-1}\left(1+\frac{p-2j}{N}\right)^2
 \le\frac{J_s^2}{\lambda_N}e^{2p/N}.
\end{equation}
The first inequality uses $4ab\le(a+b)^2$ on each positive factor
in the two binomials; sectors with $m<p$ have $C_m=0$.
Under $b_m$, $\EE_b(m-n)^2=N+1$, so
\begin{equation}
 \sum_{m\text{ central}}b_m B_m
 \le\lambda_N+\lambda_N^2+\lambda_N(1+1/N).
\end{equation}
Outside $[N/4,3N/4]$ the total binomial weight is at most
$2e^{-N/8}$. The true $F_m$ has normalized HS norm at most one.
The OU proxy also has annealed norm at most one. Hence the exterior
contribution to the squared error is at most $8e^{-N/8}$.
Use the clipped $\alpha_m$ there. Equation \eqref{eq:finite} now
proves \eqref{eq:main}. The absent charge-zero removal term may
equivalently be assigned zero error.

\paragraph{The new interface.}
The approximation replaces a neighboring-charge exponential by an
explicit conditional Gaussian average in the original sector:
\begin{equation}
 \mathcal O_\alpha(e^{ish(J)})
 =\EE_{J'}e^{is[\alpha h(J)+\sqrt{1-\alpha^2}h(J')]}.
\end{equation}
It retains every coupling chaos and the original common disorder.
It does not assert samplewise commutation with $h(J)$.

For the actual force, $a=\alpha h$ and
\begin{equation}
 \Lambda(s)=-F_m''(s)+i\{a,F_m'(s)\}+aF_m(s)a.
\end{equation}
Differentiating the \emph{OU proxy itself} is exact and gives,
with $H'=\alpha h+\sqrt{1-\alpha^2}\,u$,
\begin{equation}
 \Lambda_{\mathrm{OU}}(s)
 =(1-\alpha^2)\EE_u\big[u e^{isH'}u\big].
 \label{eq:proxyforce}
\end{equation}
The independent $u$ in \eqref{eq:proxyforce} is inside the same
exponential as its two insertions; it cannot be replaced by an
independent covariance sandwich. This is an explicit next target.

The proved error is currently unweighted $L^2$. It cannot be
differentiated to estimate $\Lambda-\Lambda_{\mathrm{OU}}$, nor
multiplied freely by $h$. Moreover, the collective factor $m$ makes
an $O(N^{-1})$ force error order one. Next we must obtain the
required derivative/energy-weighted transfer bounds and identify
the corresponding fluctuation, before discarding any such term.
The finite-time $X_0$ covariance and bulk meson dictionary remain open.

\newpage
\section{Evidence and literature boundary}
The spectrum in \eqref{eq:spectrum} is the one-body specialization
of the established embedded-GUE spectrum, not a new spectrum.
The Gaussian creation, adjoint and OU facts are standard. This
stage applies them to the concrete common-coupling adjacent-charge
map: its exact directional defects, the low-degree polar estimate,
the complete real-time propagation bound and the OU comparison.
The proofs are written above; finite examples do not establish
their uniform asymptotic scope.

\paragraph{Saved verification.}
Six systems $(N,p,m)=(3,1,2),(4,1,2),(4,2,3),(5,1,2),
(5,2,3),(5,2,4)$ were constructed directly from CAR matrices.
For each, every real Gaussian coordinate and the full matrix-input
space were retained. Complete superoperator checks verify both
channel Grams, the forward defect, both polar maps, the centered
reverse defect, its full reverse quadratic form and the complex
source--defect orthogonality. These are not Monte Carlo checks or
fits of a large-$N$ rate. Two rational identities check the forward
coefficient, one polynomial identity checks the polar amplitudes,
and 230 exact rational cases check the source variance from the
known Casimir second moment.

The three scientific JSON files are saved in
\texttt{results/susskind-memory-charged-transfer/}. Their drivers
record parameter sets, dependency versions and input hashes and
refuse to overwrite existing outputs. Small systems also include
unequal operator-space dimensions and channel zero modes. The
previous ordinary rank--chaos evidence and all older outputs remain
unchanged. Final source, PDF, two-pass build log and page renders
are bound by a separate local evidence record.

\paragraph{Sources and next checks.}
Kota's covariance result is read at finite $N$, without importing a
dilute-filling approximation. Chen--Liu is used for the explicit
Gaussian adjoint convention and Hermite structure. Standard CP
Schwarz structure was also checked against Choi's original paper;
the new rate is not obtained from a generic CP contraction alone.
The complex-SYK finite-site rules remain the complementary route
to a direct force four-point function. Their leading neutral-probe
formula cannot replace the charged, subleading calculation.

\begin{thebibliography}{9}
\bibitem{Kota} V. K. B. Kota,
\emph{SU(N) Wigner--Racah algebra for the matrix of second moments
of embedded Gaussian unitary ensemble of random matrices},
\href{https://arxiv.org/abs/nucl-th/0410069}{arXiv:nucl-th/0410069},
Eqs.~(17), (25)--(26).
\bibitem{ChenLiu} Y. Chen and Y. Liu,
\emph{On the eigenfunctions of the complex Ornstein--Uhlenbeck operators},
\href{https://arxiv.org/abs/1209.4990}{arXiv:1209.4990v3},
Sec.~2.2 and Theorem~2.6.
\bibitem{Choi} M.-D. Choi,
\emph{A Schwarz inequality for positive linear maps on $C^*$-algebras},
Illinois J. Math. \textbf{18}, 565--574 (1974),
\href{https://doi.org/10.1215/ijm/1256051007}{doi:10.1215/ijm/1256051007}.
\bibitem{BNR} M. Berkooz, V. Narovlansky and H. Raj,
\emph{Complex Sachdev--Ye--Kitaev model in the double scaling limit},
\href{https://arxiv.org/abs/2006.13983}{arXiv:2006.13983v2},
Secs.~3.2--3.3 and 7--8.
\end{thebibliography}
\end{document}
